\Article{341}
The first main task of these investigations will be to determine the number of prime functions of any given degree. For a given modulus, the number of all different (incongruent) functions of any given degree is finite. Some of these are composed of prime functions of lower degrees, while others are prime; hence the number of the latter is also finite. A rigorous development of this matter is rather delicate; we shall begin with simpler cases.

With the modulus set $=p$, the number of all different functions of degree $n$ of the form
\[
x^n+Ax^{n-1}+Bx^{n-2}+Cx^{n-3}+\text{etc.}
\]
will be $p^n$. For the number of coefficients $\pdot A,B\pdot C$, and so on, is $n$, and each can independently of the others be $\equiv0,1,2,3\dotstwo(p-1)(\mathrm{mod.}\,p)$. The theory of combinations therefore gives $p^n$ different combinations, which determine the collection of all different functions of this degree.

Thus there will be $p$ functions of degree one, namely $x,x+1,x+2$ up to $x+p-1$; there will be $pp$ functions of degree two, and so on.

\Article{342}
We have already pointed out that all functions of degree one must be regarded as prime. If, therefore, as is sufficient for our purpose, we restrict ourselves to functions whose highest term has coefficient $1$, there will be $p$ functions of the first degree, or of one dimension.

All functions of degree two are either composed of two of degree one or are prime. The theory of combinations shows that $p$ different things, with repetitions allowed, can be combined in $\frac{p\pdot p+1}{1\pdot2}$ different ways. Thus there will be that many functions composed of two prime functions of degree one, and hence $pp-\frac{p\pdot p+1}{1\pdot2}=\frac12(pp-p)$ prime functions of degree two.

\SourceLeaf{230}
\Needspace{7\baselineskip}
Similarly, from all functions of degree three, whose number is $p^3$, we must exclude those composed of three prime functions of degree one, whose number is $\frac{p\pdot p+1\pdot p+2}{1\pdot2\pdot3}$, and also those composed of a prime function of degree one and another of degree two, whose number is $p\pdot\frac12(pp-p)$. Removing these leaves $\frac13(p^3-p)$; this, then, is the number of prime functions of degree three. Clearly we can always continue in this way.

\Article{343}
To carry out these operations more easily and at the same time prepare the way for developing the general law, we shall consider the matter generally. For brevity we denote by $(1)$ the number of prime functions of degree one, by $(2)$ the number of prime functions of degree two, and further by $(1^2)$ the number of functions composed of two prime functions of degree one, and so on, and so on. In general, $(1^\alpha2^\beta3^\gamma\ldots)$ denotes the number of all functions composed of prime functions, namely $\alpha$ of degree one, $\beta$ of degree two, $\gamma$ of degree three, and so on. Their degree will therefore be $\alpha+2\beta+3\gamma+\text{etc.}$ It follows from the preceding results and the theory of combinations that
\[
(1^\alpha2^\beta3^\gamma4^\delta\ldots)=(1^\alpha)(2^\beta)(3^\gamma)(4^\delta)\dotstwo
\]
\[
(1^\alpha)=\frac{(1)\pdot(1)+1\pdot(1)+2\pdot(1)+3\ldots(1)+\alpha-1}{1\pdot2\pdot3\pdot4\ldots\alpha}
\]
or, in general,
\[
(a^\alpha)=\frac{(a)\pdot(a)+1\pdot(a)+2\pdot(a)+3\ldots(a)+\alpha-1}{1\pdot2\pdot3\pdot4\ldots\alpha}
\]

\Needspace{11\baselineskip}
Finally, if all the different ways of composing the number $n$ by addition from the numbers $1,2,3,\ldots$ are collected, and these ways are denoted by $\alpha\pdot1+\beta\pdot2+\gamma\pdot3+\text{etc.}$, then the sum of all these expressions $(1^\alpha2^\beta3^\gamma\dotstwo)$ is clearly equal to the number of all functions of degree $n$, that is, $=p^n$. Thus
\[
\begin{aligned}
p&=(1)\\
pp&=(1^2)+(2)\\
p^3&=(1^3)+(1\pdot2)+(3)\\
p^4&=(1^4)+(1^2\pdot2)+(1\pdot3)+(2^2)+(4)\\
&\text{etc.}
\end{aligned}
\]

The expression $p^n$ clearly contains, besides the quantities $(1),(2),(3)$, and so on,
\SourceLeaf{231}
also $(n)$, which shows how all these quantities are to be determined from the preceding ones. Thus we find
\[
\begin{array}{lll}
(1)=p&(4)=\frac14(p^4-pp)&(7)=\frac17(p^7-p)\\
(2)=\frac12(pp-p)&(5)=\frac15(p^5-p)&(8)=\frac18(p^8-p^4)\\
(3)=\frac13(p^3-p)&(6)=\frac16(p^6-p^3-pp+p)&\text{etc.}
\end{array}
\]

\Article{344--346}
From the beginning of this sequence we observe that the highest term of the expression $(n)$ is $\frac1n p^n$. If $n$ is prime, $-\frac1n p$ is added to it; if $n$ is composite, however, the law is less evident. But if we consider the matter more closely, we see that
\[
\begin{array}{ll}
p=(1)&p^5=5(5)+(1)\\
pp=2(2)+(1)&p^6=6(6)+3(3)+2(2)+(1)\\
p^3=3(3)+(1)&p^7=7(7)+(1)\\
p^4=4(4)+2(2)+(1)&p^8=8(8)+4(4)+2(2)+(1)\text{ etc.}
\end{array}
\]
Here the law of progression is apparent: if all the divisors of the number $n$ are $\alpha,\beta,\gamma,\delta$, and so on, then
\[
p^n=\alpha(\alpha)+\beta(\beta)+\gamma(\gamma)+\delta(\delta)+\text{etc.}
\]

We now set out to prove that this observation holds in general.

\Needspace{16\baselineskip}
We showed that the sum of all expressions of the form $(1^\alpha)(2^\beta)(3^\gamma)\ldots$, where always $\alpha+2\beta+3\gamma+\ldots=n$, exhausts all functions of degree $n$ and is therefore $=p^n$. It follows that, \textemdash\ \textemdash\ \textemdash. If
\[
\left(\frac1{1-x}\right)^{(1)}\left(\frac1{1-xx}\right)^{(2)}\left(\frac1{1-x^3}\right)^{(3)}\ldots
\]
is expanded into the series $1+Ax+Bx^2\ldots=P$, then
\[
A=p,\quad B=p^2,\quad C=p^3\text{ etc.}
\]
\[
\frac{x\,\mathrm{d}P}{P\,\mathrm{d}x}=\frac{(1)x}{1-x}+\frac{2(2)x^3}{1-x^2}+\frac{3(3)x^3}{1-x^3}\ldots
\]
\noindent\hspace*{.28\linewidth}\textemdash\ \textemdash\ \textemdash\ \textemdash

[From this, on substituting $\frac{px}{1-px}$ for $\frac{x\,\mathrm{d}P}{P\,\mathrm{d}x}$ and expanding the individual fractions into infinite series, the truth of the theorem follows of itself.]

\SourceLeaf{232}
\Article{347}
This theorem can also be expressed in another way. If all the divisors of the number $n$ are $n,1,\delta,\delta',\delta'',\delta'''$, and so on, the theorem states that
\[
p^n=n(n)+(1)+\delta(\delta)+\delta'(\delta')+\text{etc.}
\]
Now the product of the $(n)$ prime functions of degree $n$ clearly has degree $n(n)$, and likewise for the others. Hence

\emph{The product of all prime functions of degree one and of degrees $n,\delta,\delta'$, and so on, will have degree $p^n$.}

\Needspace{12\baselineskip}
It is now easy to deduce the value of the expression $(n)$ itself from this theorem, but for brevity we omit the analysis, which is not difficult. Let $n=a^\alpha b^\beta c^\gamma$, and so on, where $a,b,c$, and so on, are distinct prime numbers. Then
\[
n(n)=p^n-\Sigma p^{\frac na}+\Sigma p^{\frac n{ab}}-\Sigma p^{\frac n{abc}}\text{ etc.}
\]
Here $\Sigma p^{\frac n{abc\dotstwo}}$ denotes the collection of all expressions similar to $p^{\frac n{abc\dotstwo}}$ when the quantities $a,b,c\dotstwo$ are permuted among themselves in every way. Thus for $n=36$ we have $36(36)=p^{36}-p^{18}-p^{12}+p^6$.

Let us add one more observation. If $n$ has the form $a^\alpha$, and $a$ is prime, then $n(n)=p^n-p^{\frac na}$. Since $(n)$ must be an integer, we therefore have, whatever $p$ may be,
\[
p^n\equiv p^{\frac na}\;(\mathrm{mod.}\,n)
\]
Thus, if $p$ is coprime to $a$,
\[
p^{n-\frac na}\equiv1\;(\mathrm{mod.}\,n)
\]
and for $\alpha=1$,
\[
p^{a-1}\equiv1\;(\mathrm{mod.}\,a)
\]
It is remarkable that these theorems can be obtained in such different ways.
